\section{Resumo}

	Escoamento em meios porosos consiste em um tema de grande impacto no cenário nacional e internacional, haja vista a sua relação com o setor químico e petroquímico. O investimento no desenvolvimento de tecnologias e técnicas mais avanças para caracterização do processo de extração do petróleo é um tema recorrente, e afeta diversos setores da cadeia produtiva, sobretudo, no segmento de derivados do petróleo.
	
	No Brasil, em particular, o relatório do balanço energético nacional, publicado pela EPE (Empresa de Pesquisa Energética) mostrou que a produção nacional de petróleo cresceu 4\% em 2022, atingindo a média de 3,0 milhões de barris diários, dos quais 97\% são de origem marítima (\textit{offshore}). A parcela significativa de petróleo extraída em regiões \textit{offshore} é caracterizada por desafios que permeiam não somente a engenharia, mas também áreas correlacionadas como matemática, física, química e geofísica. A complexidade da extração em regiões que distam quilômetros da superfície do nível do mar evidencia a complexidade do processo. Assim, faz-se necessário dimensionar e demarcar possíveis campos de reservatórios de petróleo, otimizar a posição dos poços de injeção e extração, visando, sobretudo, a maximização da produção e vida útil dos reservatórios. 
	
	Com os avanços da computação científica e o advento dos supercomputadores, a simulação computacional foi incorporada como ferramenta para otimização da produção de petróleo. Escoamentos altamente complexos, que incorporam a física dos escoamentos em subsuperfícies e meios altamente heterogêneos, marcados pela presença de fraturas e carstes, passaram a ser tratados sob a ótica dos métodos numéricos e da computação de alto desempenho. 
	
	Muito embora diversos modelos para escoamento em meios porosos estejam atualmente bem caracterizados e delimitados na literatura (e.g. \textit{black-oil model} e \textit{compositional model}), todos eles possuem um elo comum, uma lei constitutiva denominada de Lei de \textit{Darcy}, a qual estabelece uma relação entre o fluxo de um determinado fluido através de um meio poroso e o gradiente de pressão neste. Esta lei constitutiva, contudo, apresenta limitações, sendo válida para escoamento em regime laminar, estacionário, isotérmico e incompressível. Este conjunto de hipóteses, evidentemente, não caracteriza todos os diversos cenários possíveis do escoamento no interior dos reservatórios de petróleo. A relação entre o fluxo e o gradiente de pressão no interior de uma rocha porosa pode apresentar características bastante complexas que se afastam, naturalmente, dessa lei constitutiva.
	
	Nesse aspecto, o presente estudo visa a ampliar um novo conceito, intitulado \textit{data-driven computing}, e introduzido por T. Kichdoerfer e M. Ortiz em 2016, no contexto do escoamento em meio porosos. Nesta abordagem, a lei constitutiva do problema é suprimida e a solução do problema é obtida diretamente dos dados experimentais, evitando assim a modelagem empírica da lei constitutiva ou lei material. Esta temática nova exibe um nítido contraste com a computação convencional, na computação baseada em dados, os pontos de dados de materiais experimentais são usados diretamente em cálculos, no lugar de um modelo de material empírico. Desta forma, o empirismo, o erro e a incerteza da modelagem de materiais são totalmente eliminados e nenhuma perda de informação experimental é incorrida. Especificamente, os solucionadores orientados por dados procuram atribuir a cada ponto material o estado de um conjunto de dados pré-especificado que está mais próximo de satisfazer as leis de conservação.
	
	Nesse viés, portanto, propõe-se a construção de soluções variacionais para o problema de escoamento em meios porosos (\textit{Darcy's flow in porous media}) utilizando este novo paradigma. O projeto de pesquisa, adicionalmente, visa estender este conceito para problemas não lineares e prolemas de deformação mecânica em regime elástico e plástico. 
	
	
	
	
	
	
	
	
	
	
	